\loai{Loại 1. Tính hiệu suất theo Q1 và A}

\loai{Loại 2. Tính hiệu xuất theo Q1 và Q2}
\begin{ex}%[1P6B4-2][Loai5] 
	Một động cơ nhiệt trả cho nguồn lạnh $70$\% nhiệt lượng thu được của nguồn nóng. Nhiệt lượng thu được trong một chu trình là $6200\ J $. Hiệu suất và công của động cơ trong một chu trình là
	\choice
	{\True 30\%; 1860 J}
	{70\%; 1860 J}
	{70\%; 4340 J}
	{30\%; 4340 J}
	\loigiai{
		\begin{itemize} 
			\item Nguồn lạnh nhận được 70\% nhiệt lượng thu được từ nguồn nóng nghĩa là "hao phí" mất 70\% hay hiệu suất là $H=30\%=0,3$.
			\item Ta có $H = \dfrac{A}{Q_N}\Leftrightarrow 0,3=\dfrac{A}{6200}\Rightarrow A=1860\ J$.
		\end{itemize}
	}
	%<MyLT2>
	\haidong{3}\end{ex}


\begin{ex}%[1P6-4.2B][Loai1] 
	Một động cơ nhiệt mỗi giây nhận từ nguồn nóng nhiệt lượng $3,2.10^4$ J đồng thời truyền cho nguồn lạnh $2,8.10^4$ J. Hiệu suất của động cơ là
	\choice
	{87,5\%}
	{50\%}
	{\True 12,5\%}
	{75\%}
	\loigiai{
		Hiệu suất động cơ: 
		$H=\dfrac{A}{Q_N}=\dfrac{\left|{Q_N-Q_L}\right|}{Q_N}=\dfrac{{3,2.10}^4-{2,8.10}^4}{{3,2.10}^4}=\dfrac{1}{8}\Rightarrow H=12,5\%$.
	}
	\haidong{5}\end{ex}

\begin{ex}%[1P6-4.2B][Loai1] 
	Một động cơ nhiệt mỗi giây nhận từ nguồn nóng nhiệt lượng $3,6. 10^4$ J đồng thời truyền cho nguồn lạnh $3,2. 10^4$ J. Hiệu suất của động cơ là
	\choice
	{\True $11\%$}
	{$1\%$}
	{$22\%$}
	{$12\%$}
	\loigiai{ Hiệu suất của động cơ: $H=\dfrac{Q_N-Q_L}{Q_N}\approx 11\%$.
	}
	\haidong{4}\end{ex}

\begin{ex} \nguon{1}
	Một động cơ nhiệt mỗi giây nhận từ nguồn nóng nhiệt lượng $4,32.10^4$ J đồng thời nhường cho nguồn lạnh $3,84.10^4\ J$. Hiệu suất của động cơ nhiệt là 
	\choice
	{$12,5 \%$}
	{$50 \%$}
	{\True $11,11 \%$}
	{$88,89 \%$}
	\loigiai{
	}
\end{ex}

\loai{Loại 3. Tính hiệu suất theo Q2 và A}

\begin{ex}
	Một động cơ nhiệt lí tưởng thực hiện một công $5 \ kJ$ đồng thời truyền cho nguồn lạnh nhiệt lượng 15 kJ. Hiệu suất của động cơ nhiệt này là
	\choice
	{$33,33 \%$}
	{$75 \%$}
	{\True $25 \%$}
	{$66,67 \%$}
	\loigiai{
		$A=Q_1-Q_2 \Rightarrow Q_1=A+Q_2=5+15=20\ kJ$
		
		$H=\dfrac{A}{Q_1}=\dfrac{5}{20}=0,25=25 \%$.
		
	}
	\haidong{5}\end{ex}

\loai{Loại 4. Từ hiệu suất, tính các đại lượng}

\begin{ex}
	Một động cơ nhiệt có hiệu suất $25 \%$, công suất $30\ kW$. Nhiệt lượng mà nó tỏa ra cho nguồn lạnh trong 5 giờ làm việc liên tục là
	\choice
	{$176.10^{7}\ J$}
	{$194 . 10^{7}\ J$}
	{$213 . 10^{7}\ J$}
	{\True $162.10^{7}\ J$}
	\loigiai{
		\begin{itemize}
			\item $A=P t=30 . 10^3 . 5 . 60 . 60=54.10^{7}\ J$.
			\item $H=\dfrac{A}{Q_1} \Rightarrow 0,25=\dfrac{54 . 10^{7}}{Q_1} \Rightarrow Q_1=216 . 10^{7}\ J$.
			\item $Q_2=Q_1-A=216 . 10^{7}-54 . 10^{7}=162 . 10^{7}\ J$.
		\end{itemize}
	}
	\haidong{4}\end{ex}



\loai{Loại 5. Tính lượng nhiên liệu tiêu thụ}

\begin{ex} \nguon{1}
	Động cơ nhiệt thực hiện công có ích $9,2.10^6 \ J$ thì phải tiêu tốn lượng xăng là $1 \ kg$. Biết khi đốt cháy hoàn toàn $1 \ kg$ xăng ta thu được nhiệt lượng $46 . 10^{6} \ J / kg$. Hiệu suất của động cơ là
	\choice
	{$15 \%$}
	{\True $20 \%$}
	{$25 \%$}
	{$30 \%$}
	\loigiai{
	}
\end{ex}
 
\begin{ex} \nguon{1}
	Động cơ nhiệt thực hiện công có ích $1,57.10^7 \ J$, phải tiêu tốn $1 \ kg$ xăng. Biết khi đốt cháy hoàn toàn $1 \ kg$ xăng ta thu được nhiệt lượng $46 . 10^{6} \ J / kg$. Hiệu suất của động cơ là
	\choice
	{$13,14 \%$}
	{$13,43 \%$}
	{\True $34,13 \%$}
	{$43,13 \%$}
	\loigiai{
	}
\end{ex}



\begin{ex} \nguon{1}
	Một ô tô có công suất $16\ kW$. Biết hiệu suất của động cơ là $20 \%$. Biết khi đốt cháy hoàn toàn $1 \ kg$ xăng ta thu được nhiệt lượng $46 . 10^{6} \ J / kg$. Khối lượng xăng tiêu hao để xe chạy trong 1 giờ là
	\choice
	{\True $6,26 \ kg$}
	{$10 \ kg$}
	{$8,2 \ kg$}
	{$20 \ kg$}
	\loigiai{
	}
\end{ex}
\loai{TLN}
\begin{ex} \nguon{1}
	Một động cơ ôtô trong thời gian 4 giờ chạy liên tục thì ôtô tiêu thụ hết 60 lít xăng. Biết hiệu suất của động cơ là $32 \%$, năng suất tỏa nhiệt của xăng là $46.10^{6} \ J / kg$ và khối
	lượng riêng của xăng là $0,7 \ kg / dm^3$. Công suất của động cơ ô tô bằng bao nhiêu kW (làm tròn kết quả đến chữ số hàng phần mười)?
	\shortans[oly]{42,9 }  
	\loigiai{
		\begin{itemize}
			\item TLN: 42,9
			\item Nhiệt lượng cung cấp khi xăng cháy hết: $Q_1=V \rho q=1932.10^{6} \ J$.
			\item Công động cơ thực hiện được:
			$
			A=Q_1 H=618,24 . 10^{6} \ J
			$
			\item Công suất của động cơ: $\mathcal{P}=\dfrac{A}{t}=42,9 . 10^3 \ W\approx 42,9 \ kW$.\end{itemize}
	}
\end{ex} 
\begin{ex}%[1P6G4-2][Loai6] 
	Một ô tô chuyển động với vận tốc $54\ km/h$ thì tiêu thụ hết $60$ lít xăng. Biết động cơ của ô tô có công suất $45\ kW$ và hiệu suất $25\%$. Năng suất toả nhiệt của xăng là $46.10^6\ J/kg$ và khối lượng riêng của xăng là $700\ kg/m^3$. Ô tô có thể đi được đoạn đường dài bao nhiêu km (làm tròn kết quả đến chữ số hàng đơn vị)?
	\shortans[oly]{ 161} 
	\loigiai{
		\begin{itemize} 
			\item Quãng đường ô tô đi được là: $s=vt$.
			\item Ta có: $t=\dfrac{A}{P}\quad (1)$.
			\item Mặt khác:  $H=\dfrac{A}{Q_1}=\dfrac{A}{mq}=\dfrac{A}{VD q} =25\% \quad (2)$.
			\item Từ $(1),\ (2)$ 
			
			$\Rightarrow t=\dfrac{25VD q}{100P}=\dfrac{25.60.10^{-3}.700.4,6.10^7}{100.45.10^3}=10733\ s\\ \Rightarrow s=vt=15.10733\approx 161000\ m=161\ km.$
		\end{itemize}
	}
	\haidong{6}\end{ex}